% Vietnamese translation for OI Public Library
% Source-PDF: 图论 GRAPH THEORY/区间图、弦图和完美图.pdf
\documentclass[11pt]{article}
\usepackage{oipl-vn}

\title{Giới thiệu đồ thị khoảng, đồ thị dây cung và đồ thị hoàn hảo}
\author{Dzx}
\date{}

\begin{document}
\maketitle

\origtitle{Introduction to interval graphs, chordal graphs and perfect graphs}
\sourcepath{Graph Theory / Interval graphs, chordal graphs and perfect graphs}

\section{Nội dung}

Bài này giới thiệu ba lớp đồ thị thường gặp trong mô hình hóa và thuật toán:
\term{đồ thị khoảng} (interval graph), \term{đồ thị dây cung} (chordal graph)
và \term{đồ thị hoàn hảo} (perfect graph). Trọng tâm là cách nhìn từ biểu diễn
hình học của các khoảng sang các tính chất tổ hợp: tô màu tối ưu, clique cực
đại, thứ tự khử hoàn hảo, nhận biết đồ thị dây cung bằng LexBFS hoặc MCS, và
nhận biết đồ thị khoảng bằng thứ tự liên tiếp của các clique cực đại.

Các hình trong slide gốc chủ yếu minh họa các khoảng trên một trục, các cạnh
giao nhau giữa chúng, và các mẫu biến đổi của PQ-tree. Ở bản ghi chú này các
hình đó được diễn giải bằng lời thay vì chèn lại ảnh; nội dung chữ của slide
được đối chiếu bằng trích xuất văn bản và các slide PQ-tree được kiểm tra
trực quan.

\section{Đồ thị khoảng và bài POJ 1083}

\subsection{Mô hình hóa bài Moving Tables}

Trong bài \textit{Moving Tables} (POJ 1083), một công ty có $400$ phòng nằm hai
bên một hành lang hẹp. Các phòng lẻ nằm ở một phía, các phòng chẵn nằm ở phía
còn lại; hai phòng đối diện nhau dùng chung cùng một đoạn hành lang. Cần chuyển
một số bàn từ phòng này sang phòng khác. Nếu hai lần chuyển bàn phải đi qua
cùng một đoạn hành lang thì chúng không thể diễn ra đồng thời. Mỗi bàn mất
$10$ phút để chuyển, hỏi thời gian ít nhất để hoàn thành tất cả các lần chuyển.

Với mỗi lần chuyển từ phòng $a$ sang phòng $b$, ta đổi số phòng về chỉ số đoạn
hành lang
\[
  p(x)=\left\lceil \frac{x}{2}\right\rceil.
\]
Lần chuyển đó chiếm các đoạn từ $\min(p(a),p(b))$ tới $\max(p(a),p(b))$, tức
một khoảng trên trục hành lang. Hai lần chuyển xung đột đúng khi hai khoảng
tương ứng giao nhau. Vì vậy ta dựng một đồ thị: mỗi lần chuyển là một đỉnh, hai
đỉnh nối cạnh nếu hai khoảng của chúng giao nhau. Cần chia các đỉnh thành ít
nhất số nhóm sao cho trong mỗi nhóm không có cạnh; đó chính là bài toán tô màu
đồ thị.

Với đồ thị tổng quát, tìm số màu nhỏ nhất là NP-khó. Nhưng đồ thị vừa dựng là
đồ thị khoảng, nên số màu nhỏ nhất bằng kích thước clique lớn nhất, tức bằng
số khoảng lớn nhất cùng phủ một đoạn hành lang. Đáp án của bài là
\[
  10\cdot \max_x \#\{\text{lần chuyển đi qua đoạn }x\}.
\]

\subsection{Định nghĩa}

Một \term{khoảng} là một đoạn trên trục số thực có hai đầu mút; đầu mút có thể
mở hoặc đóng. Cho một họ khoảng $\{I_v\}$, ta định nghĩa \term{đồ thị giao}
của họ này như sau: mỗi khoảng $I_v$ ứng với một đỉnh $v$, và giữa hai đỉnh
$v,w$ có cạnh khi và chỉ khi
\[
  I_v\cap I_w\ne \varnothing.
\]

Một đồ thị $G$ được gọi là \term{đồ thị khoảng} nếu nó là đồ thị giao của một
họ khoảng nào đó.

Các slide gốc đưa ra hai hình đối chiếu: một họ khoảng có thể biểu diễn đúng
một đồ thị khoảng, và một đồ thị khác không thể biểu diễn bằng các khoảng trên
cùng một trục. Điểm quan trọng là quan hệ kề phải đến từ giao nhau trên một
đường thẳng, không phải từ một cách vẽ tùy ý của đồ thị.

\subsection{Các dạng đầu mút và thứ tự tự nhiên}

\paragraph{Định lý.}
Các đồ thị sinh bởi khoảng mở, khoảng đóng và khoảng nửa mở nửa đóng là cùng
một lớp đồ thị.

\paragraph{Ý tưởng chứng minh.}
Vì các khoảng nằm trên trục số thực liên tục, ta có thể dịch hoặc co giãn các
đầu mút một lượng rất nhỏ để loại bỏ trường hợp giao nhau chỉ tại đầu mút, mà
không làm thay đổi các quan hệ giao nhau cần giữ. Do đó kiểu mở hay đóng của
đầu mút không làm thay đổi lớp đồ thị thu được.

\paragraph{Hệ quả.}
Mọi đồ thị khoảng đều có một biểu diễn bằng khoảng trong đó không có hai đầu
mút nào trùng nhau. Khi đó ta có thể sắp xếp các đỉnh theo thứ tự tăng dần của
đầu mút trái của khoảng đại diện. Thứ tự này được gọi là \term{thứ tự tự nhiên}
của đồ thị khoảng.

Ký hiệu thứ tự tự nhiên là
\[
  v_1,v_2,\ldots,v_n.
\]
Với một đỉnh $v_i$, đặt
\[
  \operatorname{Pred}(v_i)=
  \{v_j\mid (v_i,v_j)\in E,\ j<i\}.
\]
Ta gọi đây là tập các \term{tiền nhiệm} của $v_i$ trong thứ tự đang xét.

\paragraph{Tính chất then chốt.}
Với thứ tự tự nhiên của một đồ thị khoảng,
\[
  \{v_i\}\cup \operatorname{Pred}(v_i)
\]
là một clique. Thật vậy, mọi tiền nhiệm của $v_i$ có đầu mút trái nằm trước
đầu mút trái của $v_i$ nhưng vẫn giao với $I_{v_i}$, nên chúng đều chứa điểm
đầu trái của $I_{v_i}$. Do đó chúng đôi một giao nhau.

\section{Tô màu và clique lớn nhất trên đồ thị khoảng}

Từ thứ tự tự nhiên $v_1,\ldots,v_n$, ta tô màu tham lam:
\begin{enumerate}
  \item Duyệt các đỉnh theo thứ tự $v_1,v_2,\ldots,v_n$.
  \item Khi xét $v_i$, gán cho nó màu nhỏ nhất chưa xuất hiện trong
  $\operatorname{Pred}(v_i)$.
\end{enumerate}

Thuật toán này cho một tô màu tối ưu của đồ thị khoảng. Lý do là tại thời điểm
xét $v_i$, tập $\{v_i\}\cup \operatorname{Pred}(v_i)$ là một clique, nên nếu
$\operatorname{Pred}(v_i)$ đã dùng $k$ màu khác nhau thì mọi cách tô màu đều
cần ít nhất $k+1$ màu cho clique đó. Thuật toán tham lam dùng đúng số màu cần
thiết ở từng bước, nên tổng số màu nó dùng bằng kích thước clique lớn nhất:
\[
  \chi(G)=\omega(G).
\]
Ở đây $\chi(G)$ là số màu nhỏ nhất, còn $\omega(G)$ là kích thước clique lớn
nhất của $G$.

Với bài POJ 1083, clique lớn nhất chính là số lần chuyển cùng đi qua một đoạn
hành lang. Vì vậy không cần giải bài tô màu tổng quát; chỉ cần đếm độ phủ lớn
nhất của các khoảng.

\section{Thứ tự khử hoàn hảo}

\term{Thứ tự khử hoàn hảo} (perfect elimination order, PEO) là một thứ tự
$v_1,\ldots,v_n$ sao cho với mọi $i$, tập $\operatorname{Pred}(v_i)$ là một
clique. Theo tính chất ở trên, thứ tự tự nhiên của mọi đồ thị khoảng là một
thứ tự khử hoàn hảo.

PEO là công cụ trung tâm để xử lý các bài toán tối ưu trên đồ thị dây cung.
Khi đã có PEO, nhiều bài toán vốn khó trên đồ thị tổng quát trở nên đơn giản
hơn, chẳng hạn:
\begin{itemize}
  \item tìm clique lớn nhất;
  \item tô màu tối ưu;
  \item tìm tập độc lập lớn nhất;
  \item phủ đỉnh nhỏ nhất trong một số lớp bài toán liên quan;
  \item phủ bằng số clique ít nhất.
\end{itemize}

Với quy ước PEO dùng tập tiền nhiệm như trên, tô màu tham lam theo thứ tự PEO
cũng tối ưu trên đồ thị dây cung. Nếu dùng quy ước phổ biến khác là ``các láng
giềng xuất hiện sau tạo thành clique'', thì thứ tự đó chính là thứ tự ngược
lại của quy ước trong bài này.

\section{Đồ thị dây cung}

\subsection{Định nghĩa}

Một đồ thị được gọi là \term{đồ thị dây cung} nếu nó không có chu trình cảm
sinh độ dài ít nhất $4$. Nói tương đương, mọi chu trình đơn độ dài ít nhất $4$
đều có một \term{dây cung}, tức một cạnh nối hai đỉnh không liên tiếp trên chu
trình.

Mọi đồ thị khoảng đều là đồ thị dây cung, vì thứ tự tự nhiên của nó là PEO, và
như định lý dưới đây cho thấy, có PEO là đúng đặc trưng của đồ thị dây cung.

\subsection{Liên hệ với thứ tự khử hoàn hảo}

\paragraph{Định lý.}
Nếu một đồ thị $G$ có thứ tự khử hoàn hảo thì $G$ là đồ thị dây cung.

\paragraph{Chứng minh phác thảo.}
Giả sử tồn tại một chu trình cảm sinh $C$ độ dài ít nhất $4$. Lấy đỉnh $v_i$
có chỉ số lớn nhất trên chu trình. Hai láng giềng của $v_i$ trên $C$ đều nằm
trước $v_i$, nên theo PEO chúng phải kề nhau. Khi đó ta có một dây cung của
$C$, mâu thuẫn với việc $C$ là chu trình cảm sinh.

\paragraph{Định lý.}
Một đồ thị $G$ là đồ thị dây cung khi và chỉ khi $G$ có thứ tự khử hoàn hảo.

Chiều còn lại dựa trên khái niệm điểm đơn thuần. Một đỉnh $v$ được gọi là
\term{điểm đơn thuần} nếu mọi láng giềng của $v$ tạo thành một clique.

\paragraph{Bổ đề.}
Mọi đồ thị dây cung đều có ít nhất một điểm đơn thuần. Nếu đồ thị dây cung đó
không phải đồ thị đầy đủ, nó có ít nhất hai điểm đơn thuần không kề nhau.

\paragraph{Bổ đề.}
Mọi đồ thị con cảm sinh của một đồ thị dây cung vẫn là đồ thị dây cung.

Từ hai bổ đề này, ta xây dựng PEO bằng cách lặp lại thao tác chọn một điểm đơn
thuần, đặt nó vào cuối thứ tự khử theo quy ước ``láng giềng phía sau tạo thành
clique'', rồi xóa nó khỏi đồ thị. Đảo ngược thứ tự thu được sẽ phù hợp với quy
ước $\operatorname{Pred}$ của bài này.

\section{Nhận biết đồ thị dây cung}

Có hai thuật toán tuyến tính thường dùng để sinh một thứ tự ứng viên cho PEO:
\term{MCS} (maximum cardinality search) và \term{LexBFS} (lexicographical
breadth-first search). Sau đó ta kiểm tra thứ tự này trong $O(n+m)$.

\subsection{Maximum cardinality search}

MCS duy trì cho mỗi đỉnh chưa chọn một trọng số, ban đầu bằng $0$. Lặp lại
$n$ lần:
\begin{enumerate}
  \item chọn một đỉnh chưa chọn có trọng số lớn nhất;
  \item đưa nó vào thứ tự;
  \item tăng trọng số của mọi láng giềng chưa chọn của nó lên $1$.
\end{enumerate}

Trọng số của một đỉnh chính là số láng giềng đã được chọn trước đó. Với đồ thị
dây cung, thứ tự do MCS sinh ra là một PEO sau khi chọn đúng chiều thứ tự.

\subsection{LexBFS}

LexBFS cũng chọn từng đỉnh một, nhưng nhãn của mỗi đỉnh là một dãy số thay vì
một trọng số đơn. Ban đầu mọi nhãn rỗng. Ở bước $i$, chọn một đỉnh chưa chọn
có nhãn lớn nhất theo thứ tự từ điển, rồi thêm chỉ số $i$ vào nhãn của mọi
láng giềng chưa chọn của nó. Slide gốc minh họa một lần chạy LexBFS cho thứ
tự
\[
  A,D,C,B,E,F,G,H,J,K,I,L.
\]

Để đạt $O(n+m)$, các đỉnh chưa chọn được chia thành các \term{xô} theo nhãn.
Một xô chứa các đỉnh có cùng nhãn; danh sách xô được giữ theo thứ tự giảm dần
của nhãn.
\begin{enumerate}
  \item Lấy đỉnh đầu tiên $v_i$ trong xô đầu tiên; đây là một đỉnh có nhãn lớn
  nhất hiện tại.
  \item Xóa $v_i$ khỏi xô của nó; nếu xô rỗng thì xóa xô.
  \item Với mỗi láng giềng chưa chọn $w$ của $v_i$, chuyển $w$ từ xô mang
  nhãn $L(w)$ sang xô mang nhãn $L(w)\circ i$.
  \item Nếu xô mới chưa tồn tại ngay trước xô cũ trong danh sách, tạo xô mới
  tại vị trí đó.
  \item Nếu xô cũ trở nên rỗng, xóa nó.
\end{enumerate}
Ở đây $L(w)\circ i$ là nhãn thu được bằng cách nối thêm $i$ vào nhãn hiện tại.

\subsection{Kiểm tra một thứ tự}

Sau khi MCS hoặc LexBFS sinh ra một thứ tự $v_1,\ldots,v_n$, ta cần kiểm tra
nó có phải PEO hay không. Với quy ước tiền nhiệm, với mỗi $v_i$ đặt $p(v_i)$
là tiền nhiệm có chỉ số lớn nhất của $v_i$ nếu tồn tại. Khi đó thứ tự là PEO
nếu và chỉ nếu với mọi tiền nhiệm $u$ khác $p(v_i)$ của $v_i$, ta có cạnh
$(u,p(v_i))$.

Điều kiện này kiểm tra rằng toàn bộ $\operatorname{Pred}(v_i)$ là clique,
nhưng chỉ cần so sánh với một đỉnh đại diện $p(v_i)$, vì nếu thứ tự là PEO thì
các kiểm tra tương tự ở các bước trước đã bảo đảm phần còn lại. Có thể cài đặt
kiểm tra trong $O(n+m)$ bằng danh sách kề và mảng đánh dấu tạm thời.

Ứng dụng trực tiếp là bài \textit{Fishing Net} trên ZOJ: cho một đồ thị, hãy
quyết định nó có phải đồ thị dây cung hay không. Ta chạy MCS hoặc LexBFS, rồi
kiểm tra PEO.

\section{Nhận biết đồ thị khoảng}

\subsection{Ba đặc trưng tương đương}

Với một đồ thị $G$, các mệnh đề sau là tương đương:
\begin{enumerate}
  \item $G$ là đồ thị khoảng.
  \item $G$ là đồ thị dây cung và $G$ là đồ thị co-comparability, nghĩa là đồ
  thị bù $\overline G$ là một đồ thị comparability.
  \item Các clique cực đại của $G$ có thể được đánh số liên tiếp
  $C_1,C_2,\ldots,C_k$ sao cho với mọi đỉnh $v$, tập chỉ số
  \[
    \{j\mid v\in C_j\}
  \]
  là một đoạn các số nguyên liên tiếp.
\end{enumerate}

Một đồ thị vô hướng được gọi là \term{comparability graph} nếu các cạnh của nó
có thể được định hướng sao cho hướng thu được không có chu trình và có tính
bắc cầu: nếu có cung $x\to y$ và $y\to z$ thì cũng có cung $x\to z$.

Hàm ý $(3)\Rightarrow(1)$ rất trực tiếp. Nếu các clique cực đại đã được xếp
theo thứ tự liên tiếp, ta gán cho mỗi đỉnh $v$ khoảng
\[
  I(v)=
  \left[
    \min\{i\mid v\in C_i\},
    \max\{i\mid v\in C_i\}
  \right].
\]
Hai đỉnh kề nhau khi và chỉ khi chúng cùng nằm trong một clique cực đại nào
đó, tức hai khoảng tương ứng giao nhau.

Hàm ý $(1)\Rightarrow(2)$ đến từ việc đồ thị khoảng là đồ thị dây cung, còn
đồ thị bù của nó có thể định hướng theo thứ tự tương đối của các khoảng không
giao nhau: nếu $I_x$ hoàn toàn nằm bên trái $I_y$, định hướng cạnh bù thành
$x\to y$. Hướng này là không chu trình và bắc cầu.

Với hàm ý $(2)\Rightarrow(3)$, slide gốc xây dựng một đồ thị có hướng $H$ trên
tập các clique cực đại. Lấy $\overline G$ sau khi đã định hướng bắc cầu và
không chu trình thành $G'$. Đặt $V(H)$ là tập các clique cực đại của $G$; với
hai clique $C_1,C_2$, thêm cung $C_1\to C_2$ nếu tồn tại $x\in C_1$,
$y\in C_2$ sao cho $x\to y$ trong $G'$. Các định lý trên slide khẳng định:
$H$ có tính bắc cầu, $H$ không có chu trình, và một thứ tự topo
$C_1,C_2,\ldots,C_k$ của $H$ chính là thứ tự clique thỏa điều kiện liên tiếp.

\subsection{Liệt kê clique cực đại trong đồ thị dây cung}

Giả sử $G$ là đồ thị dây cung và ta đã có một PEO
$v_1,\ldots,v_n$ theo quy ước tiền nhiệm. Khi đó mỗi clique cực đại của $G$ có
dạng
\[
  \{v_i\}\cup \operatorname{Pred}(v_i)
\]
với một chỉ số $i$ nào đó.

Không phải mọi tập ở dạng trên đều là clique cực đại. Nó là clique cực đại khi
và chỉ khi không có đỉnh xuất hiện sau nào chứa toàn bộ tập này trong clique
của nó. Cụ thể, với mọi hậu nhiệm $v_j$ của $v_i$ kề với $v_i$, phải tồn tại
một tiền nhiệm của $v_i$ không phải là tiền nhiệm của $v_j$. Nếu ngược lại mọi
tiền nhiệm của $v_i$ cũng kề với $v_j$, thì clique của $v_i$ bị chứa trong
clique lớn hơn ứng với $v_j$.

Nhờ vậy, sau khi nhận biết chordal và có PEO, ta có thể liệt kê các clique cực
đại để chuyển sang bước kiểm tra tính liên tiếp.

\subsection{Tính chất các số 1 liên tiếp}

\term{Tính chất các số 1 liên tiếp} (consecutive ones property, C1P hoặc COP)
của một ma trận nhị phân nói rằng có thể hoán vị các cột sao cho trong mỗi
hàng, các phần tử bằng $1$ xuất hiện thành một khối liên tiếp.

Với bài toán nhận biết đồ thị khoảng, lập ma trận $A$ có hàng là các đỉnh và
cột là các clique cực đại:
\[
  a_{ij}=1 \Longleftrightarrow v_i\in C_j.
\]
Khi đó đồ thị dây cung $G$ là đồ thị khoảng khi và chỉ khi ma trận này có C1P.
Slide gốc nhắc tới POJ 2790 như một bài kiểm tra trực tiếp xem một ma trận có
tính chất C1P hay không.

Ví dụ, nếu $N=4$ và ràng buộc đầu tiên là tập $S_1=\{2,3\}$ phải đứng liên
tiếp, các hoán vị hợp lệ của $1,2,3,4$ là những hoán vị trong đó $2$ và $3$
đứng cạnh nhau. Nếu thêm ràng buộc $S_2=\{3,4\}$, tập hoán vị còn lại bị thu
hẹp hơn nữa, chẳng hạn các thứ tự như
\[
  \langle 1,2,3,4\rangle,\quad
  \langle 1,4,3,2\rangle,\quad
  \langle 2,3,4,1\rangle,\quad
  \langle 4,3,2,1\rangle
\]
đều thỏa cả hai ràng buộc. Đây chính là kiểu thông tin mà PQ-tree lưu trữ một
cách nén.

\subsection{PQ-tree}

PQ-tree là một cây biểu diễn tập tất cả các hoán vị cột còn có thể thỏa các
ràng buộc C1P đã xử lý. Lá của cây là các cột. Nút trong có hai loại:
\begin{itemize}
  \item \term{P-node}: các con của nó có thể được hoán vị tùy ý;
  \item \term{Q-node}: thứ tự các con của nó cố định tới phép đảo ngược.
\end{itemize}

Khi thêm một ràng buộc mới $S$ yêu cầu các lá thuộc $S$ phải liên tiếp, thuật
toán đánh dấu các nút theo quan hệ với $S$:
\begin{itemize}
  \item \term{full}: toàn bộ lá trong cây con thuộc $S$;
  \item \term{empty}: không có lá nào trong cây con thuộc $S$;
  \item \term{partial}: chỉ một phần lá trong cây con thuộc $S$.
\end{itemize}
Nút thấp nhất bao trùm toàn bộ phần liên quan tới $S$ được gọi là
\term{pertinent root}. Từ nút này, cây được rút gọn bằng một số mẫu thay thế
cục bộ.

Các slide gốc liệt kê các mẫu chuẩn:
\begin{itemize}
  \item L1: lá thuộc $S$ được đánh dấu full;
  \item P1: P-node có toàn bộ con full thì trở thành full;
  \item P2: P-node là pertinent root với các con full và empty; gom các con
  full dưới một P-node mới nếu có nhiều hơn một con full;
  \item P3: P-node không phải pertinent root với con full và empty; biến nó
  thành Q-node partial, gom nhóm full và nhóm empty bằng các P-node phụ khi
  cần;
  \item P4, P5, P6: các trường hợp P-node có một hoặc hai con partial, cộng
  thêm nhóm full và empty, tùy nó có phải pertinent root hay không;
  \item Q1: Q-node có toàn bộ con full thì trở thành full;
  \item Q2: Q-node có không quá một con partial, một khối full liên tiếp, và
  phần còn lại empty;
  \item Q3: Q-node có hai con partial ở hai đầu của một khối full liên tiếp,
  phần còn lại empty.
\end{itemize}

Nếu một ràng buộc không khớp được với mẫu hợp lệ nào, ma trận không có C1P và
đồ thị không phải đồ thị khoảng. Nếu xử lý hết các hàng, bất kỳ thứ tự lá nào
do PQ-tree còn lại biểu diễn đều cho một thứ tự các clique cực đại thỏa điều
kiện liên tiếp, từ đó dựng được biểu diễn khoảng.

\section{Đồ thị hoàn hảo}

Một đồ thị $G$ được gọi là \term{đồ thị hoàn hảo} nếu với mọi đồ thị con cảm
sinh $H$ của $G$, ta đều có
\[
  \omega(H)=\chi(H).
\]
Ở đây $\omega(H)$ là kích thước clique lớn nhất của $H$, còn $\chi(H)$ là số
màu nhỏ nhất của $H$.

Đồ thị khoảng là hoàn hảo vì mọi đồ thị con cảm sinh của đồ thị khoảng vẫn là
đồ thị khoảng, và trên đồ thị khoảng ta có $\omega=\chi$. Rộng hơn, đồ thị dây
cung cũng là đồ thị hoàn hảo: mọi đồ thị con cảm sinh của nó vẫn dây cung, có
PEO, và tô màu tham lam theo PEO đạt đúng kích thước clique lớn nhất.

\paragraph{Định lý đồ thị hoàn hảo mạnh.}
Một đồ thị là hoàn hảo khi và chỉ khi nó không chứa chu trình lẻ cảm sinh độ
dài ít nhất $5$ và cũng không chứa phần bù của một chu trình như vậy dưới dạng
đồ thị con cảm sinh. Các chu trình lẻ cảm sinh độ dài ít nhất $5$ thường được
gọi là \term{odd holes}; phần bù của chúng là \term{odd antiholes}.

Định lý này đặt đồ thị khoảng và đồ thị dây cung vào bức tranh rộng hơn: chúng
là các lớp con có cấu trúc mạnh của lớp đồ thị hoàn hảo, nơi nhiều bài toán
tối ưu hóa trên đồ thị có thể được giải hiệu quả nhờ các đặc trưng thứ tự và
clique.

\end{document}
